\documentclass{amsart}
% For dealing with references we use the comment environment
\usepackage{verbatim}
\newenvironment{reference}{\comment}{\endcomment}
%\newenvironment{reference}{}{}
\newenvironment{slogan}{\comment}{\endcomment}
\newenvironment{history}{\comment}{\endcomment}

% For commutative diagrams we use Xy-pic
\usepackage[all]{xy}

% We use 2cell for 2-commutative diagrams.
\xyoption{2cell}
\UseAllTwocells

% We use multicol for the list of chapters between chapters
\usepackage{multicol}

% This is generally recommended for better output
\usepackage{lmodern}
\usepackage[T1]{fontenc}

% For cross-file-references

% Package for hypertext links:
\usepackage{hyperref}

% For any local file, say "hello.tex" you want to link to please

% Theorem environments.
%
\theoremstyle{plain}
\newtheorem{theorem}[subsection]{Theorem}
\newtheorem{proposition}[subsection]{Proposition}
\newtheorem{lemma}[subsection]{Lemma}

\theoremstyle{definition}
\newtheorem{definition}[subsection]{Definition}
\newtheorem{example}[subsection]{Example}
\newtheorem{exercise}[subsection]{Exercise}
\newtheorem{situation}[subsection]{Situation}

\theoremstyle{remark}
\newtheorem{remark}[subsection]{Remark}
\newtheorem{remarks}[subsection]{Remarks}

\numberwithin{equation}{subsection}

% Macros
%
\def\lim{\mathop{\mathrm{lim}}\nolimits}
\def\colim{\mathop{\mathrm{colim}}\nolimits}
\def\Spec{\mathop{\mathrm{Spec}}}
\def\Hom{\mathop{\mathrm{Hom}}\nolimits}
\def\Ext{\mathop{\mathrm{Ext}}\nolimits}
\def\SheafHom{\mathop{\mathcal{H}\!\mathit{om}}\nolimits}
\def\SheafExt{\mathop{\mathcal{E}\!\mathit{xt}}\nolimits}
\def\Sch{\mathit{Sch}}
\def\Mor{\mathop{\mathrm{Mor}}\nolimits}
\def\Ob{\mathop{\mathrm{Ob}}\nolimits}
\def\Sh{\mathop{\mathit{Sh}}\nolimits}
\def\NL{\mathop{N\!L}\nolimits}
\def\CH{\mathop{\mathrm{CH}}\nolimits}
\def\proetale{{pro\text{-}\acute{e}tale}}
\def\etale{{\acute{e}tale}}
\def\QCoh{\mathit{QCoh}}
\def\Ker{\mathop{\mathrm{Ker}}}
\def\Im{\mathop{\mathrm{Im}}}
\def\Coker{\mathop{\mathrm{Coker}}}
\def\Coim{\mathop{\mathrm{Coim}}}

% Boxtimes
%
\DeclareMathSymbol{\boxtimes}{\mathbin}{AMSa}{"02}

%
% Macros for moduli stacks/spaces
%
\def\QCohstack{\mathcal{QC}\!\mathit{oh}}
\def\Cohstack{\mathcal{C}\!\mathit{oh}}
\def\Spacesstack{\mathcal{S}\!\mathit{paces}}
\def\Quotfunctor{\mathrm{Quot}}
\def\Hilbfunctor{\mathrm{Hilb}}
\def\Curvesstack{\mathcal{C}\!\mathit{urves}}
\def\Polarizedstack{\mathcal{P}\!\mathit{olarized}}
\def\Complexesstack{\mathcal{C}\!\mathit{omplexes}}
% \Pic is the operator that assigns to X its picard group, usage \Pic(X)
% \Picardstack_{X/B} denotes the Picard stack of X over B
% \Picardfunctor_{X/B} denotes the Picard functor of X over B
\def\Pic{\mathop{\mathrm{Pic}}\nolimits}
\def\Picardstack{\mathcal{P}\!\mathit{ic}}
\def\Picardfunctor{\mathrm{Pic}}
\def\Deformationcategory{\mathcal{D}\!\mathit{ef}}

\setlength{\emergencystretch}{3em}
\begin{document}
\title[Stacks Project, Tag 0EC0]{357 --- Kunz flat-Frobenius criterion for regularity}
\maketitle
\noindent Copyright (C) 2005--2025 Johan de Jong. Stacks Project authors. Source: \href{https://stacks.math.columbia.edu/tag/0EC0}{Tag 0EC0}.

\noindent Editorial difficulty note (not author text). Difficulty H3: The converse transports independence of minimal generators through the flat Frobenius map. It computes the length of the Frobenius parameter quotient and compares the Hilbert-Samuel polynomial at scaled powers to force embedding dimension equal to dimension. The length asymptotics and independence argument are the decisive nonroutine obligations..

\noindent Editorial provenance note (not author text). History: excerpt from revision \texttt{a0444\allowbreak 6e57e\allowbreak c1fbc\allowbreak 252a8\allowbreak 71afc\allowbreak ec775\allowbreak 2fb28\allowbreak 07b14}, local-cohomology.tex, lines 3721--3776. Reference presentation was adapted; original-author context and core support blocks were retained or added. The mathematical wording is unchanged. Licensed under GNU FDL 1.2 or later; no invariant sections or cover texts. See ../licenses/COPYING. Context blocks reproduce author definitions and supporting results; they are not additional counted problems.

\section{Frobenius action}
\label{section-frobenius}

\noindent
Let $p$ be a prime number. Let $A$ be a ring with $p = 0$ in $A$.
The {\it Frobenius endomorphism} of $A$ is the map
$$
F : A \longrightarrow A,
\quad
a \longmapsto a^p
$$
In this section we prove lemmas on modules which have
Frobenius actions.



\noindent
In this section, we say elements $f_1, \ldots, f_r$ of a ring $A$
are {\it independent} if $\sum a_if_i = 0$ implies
$a_i \in (f_1, \ldots, f_r)$. In other words, with $I = (f_1, \ldots, f_r)$
we have $I/I^2$ is free over $A/I$ with basis $f_1, \ldots, f_r$.


\begin{lemma}[Kunz]
\label{lemma-frobenius-flat-regular}
\begin{reference}
\href{https://stacks.math.columbia.edu/bibliography}{Bibliography: Kunz-flat}
\end{reference}
Let $p$ be a prime number.
Let $A$ be a Noetherian ring with $p = 0$.
The following are equivalent
\begin{enumerate}
\item $A$ is regular, and
\item $F : A \to A$, $a \mapsto a^p$ is flat.
\end{enumerate}
\end{lemma}

\begin{proof}
Observe that $\Spec(F) : \Spec(A) \to \Spec(A)$ is the identity map.
Being regular is defined in terms of the local rings and being flat
is something about local rings, see
Algebra, Lemma \href{https://stacks.math.columbia.edu/tag/00HT}{lemma flat localization (Tag 00HT)}.
Thus we may and do assume $A$ is a Noetherian
local ring with maximal ideal $\mathfrak m$.

\medskip\noindent
Assume $A$ is regular. Let $x_1, \ldots, x_d$ be a
system of parameters for $A$. Applying $F$ we find
$F(x_1), \ldots, F(x_d) = x_1^p, \ldots, x_d^p$,
which is a system of parameters for $A$. Hence $F$ is flat, see
Algebra, Lemmas \href{https://stacks.math.columbia.edu/tag/00R4}{lemma CM over regular flat (Tag 00R4)} and
\href{https://stacks.math.columbia.edu/tag/00NQ}{lemma regular ring CM (Tag 00NQ)}.

\medskip\noindent
Conversely, assume $F$ is flat. Write $\mathfrak m = (x_1, \ldots, x_r)$
with $r$ minimal. Then $x_1, \ldots, x_r$ are independent in the sense
defined above. Since $F$ is flat, we see that $x_1^p, \ldots, x_r^p$
are independent, see Lemma \href{https://stacks.math.columbia.edu/tag/0EBZ}{lemma flat extension independent (Tag 0EBZ)}.
Hence $\text{length}_A(A/(x_1^p, \ldots, x_r^p)) = p^r$ by
Lemma \href{https://stacks.math.columbia.edu/tag/0EBY}{lemma 3 (Tag 0EBY)}.
Let $\chi(n) = \text{length}_A(A/\mathfrak m^n)$ and recall
that this is a numerical polynomial of degree $\dim(A)$, see
Algebra, Proposition \href{https://stacks.math.columbia.edu/tag/00KQ}{proposition dimension (Tag 00KQ)}.
Choose $n \gg 0$. Observe that
$$
\mathfrak m^{pn + pr} \subset F(\mathfrak m^n)A \subset \mathfrak m^{pn}
$$
as can be seen by looking at monomials in $x_1, \ldots, x_r$. We have
$$
A/F(\mathfrak m^n)A = A/\mathfrak m^n \otimes_{A, F} A
$$
By flatness of $F$ this has length $\chi(n) \text{length}_A(A/F(\mathfrak m)A)$
(Algebra, Lemma \href{https://stacks.math.columbia.edu/tag/02M1}{lemma pullback module (Tag 02M1)})
which is equal to $p^r\chi(n)$ by the above. We conclude
$$
\chi(pn + pr) \geq p^r\chi(n) \geq \chi(pn)
$$
Looking at the leading terms this implies $r = \dim(A)$, i.e., $A$ is regular.
\end{proof}
\end{document}
